% Source: https://artofproblemsolving.com/community/c3h1781263_pattern_check
% Stable identity: aops-topic-1781263-q01
\textbf{The pattern.} For $k=0,1,2,\dots$ let
\[
a=1\underbrace{6\cdots6}_{k},\qquad b=5\underbrace{0\cdots0}_{k},\qquad c=\underbrace{3\cdots3}_{k+1}.
\]
Then $a^3+b^3+c^3$ is the number obtained by writing $a$, $b$, $c$ one after another; for $k=3$, for instance, $1666^3+5000^3+3333^3=166650003333$.

\textbf{Proof.} Let $x=10^{k+1}$. Then
\[
a=\dfrac{x-4}{6},\qquad b=\dfrac {x}{2},\qquad c=\dfrac{x-1}{3}
\]
(for example $\dfrac{1000-4}{6}=166$, $\dfrac{1000}{2}=500$, $\dfrac{999}{3}=333$). Since $b$ and $c$ both have exactly $k+1$ digits, writing $a,b,c$ one after another gives the number $ax^2+bx+c$. Now
\[
216\left(a^3+b^3+c^3\right)=(x-4)^3+27x^3+8(x-1)^3=36x^3-36x^2+72x-72,
\]
so $a^3+b^3+c^3=\dfrac{x^3-x^2+2x-2}{6}$. On the other hand
\[
ax^2+bx+c=\dfrac{(x-4)x^2}{6}+\dfrac{x^2}{2}+\dfrac{x-1}{3}=\dfrac{x^3-4x^2+3x^2+2x-2}{6}=\dfrac{x^3-x^2+2x-2}{6}.
\]
The two expressions are equal, which proves the pattern for every $k\ge0$. $\blacksquare$
